Although it looks that transformers are a model that should be able to perform linear transformations in a native way, it is not the case.
Applying a linear transformation to a vector is not direct since in every step of $\ATTN$ and $\FF$ the result of the layer is added to the original vector instead of replacing it.
An intuitive idea would be to apply the transformation $M-I$ when we want to transform the vector $x$ to $Mx$, since $x + (M-I)x = x + Mx - Ix = Mx$.
This does not hold in our model since the addition and multiplication is not commutative in $\Floating$.
For avoiding the representation error, we will extend the embedding dimension to get more memory per vector such that we can store in new coordinates the value of $Mx$ and then replace it in the original ones.

\bigskip

Let $f$ be the number of a vector. \rafa{uso f porque viene de flagged}
We show a mechanism to replace it with the result of applying the linear transformation $M$ to it.
We perform the operation with two consecutive $\ATTN$ layers where $h_f$ only attends to itself.
Therefore, after these two layers $h_f$ will be replaced will have value $Mh_f$.

We need to double the embedding size from $d$ to $2d$ (plus some flags) since we need to store the result of the transformation in some coordinates (the added $d+1$ to $2d$) as a middle step.
In the first $\ATTN$ layer, we compute the transformation saving its result in the coordinates $d+1$ to $2d$.
Then, in the second layer we move the result to the right place.

\[
    \widetilde{h_f} = \left( \begin{matrix} (h_f)_1    \\ \vdots \\ (h_f)_d    \\ 0      \\ \vdots \\ 0       \end{matrix} \right) 
    \xrightarrow{\text{first $\ATTN$ layer}}
    \widetilde{h_f} = \left( \begin{matrix} 0      \\ \vdots \\ 0      \\ (Mh_f)_1 \\ \vdots \\ (Mh_f)_d \end{matrix} \right)
    \xrightarrow{\text{second $\ATTN$ layer}}
    \widetilde{h_f} = \left( \begin{matrix} (Mh_f)_1 \\ \vdots \\ (Mh_f)_d \\ 0      \\ \vdots \\ 0      \end{matrix} \right)
\]

We assume that the $d+1$ to $2d$ coordinates (used as memory) are set to $0$ by the layers previous to this mechanism.
This can be achieved maintaining them as $0$ since the beginning or transformed to $0$ with the primitive described in the section \ref{const-fun}.

\medskip

Let us construct each of the $\ATTN$ layers.
As said before, in both layers we will need for $\widetilde{h_f}$ to only attend to itself.
Therefore, we define $W_Q$ ans $W_K$ using the strategy of the section \ref{ignore-vecs}.


% %%%%%%%%%%%%%%%%%%%%%%%%%%%%%%%%%%%%%%%%%%%%%%%%%%%%
% %%%%%%%%%%% START: ONLY ATTEND TO ITSELF %%%%%%%%%%%
% %%%%%%%%%%%%%%%%%%%%%%%%%%%%%%%%%%%%%%%%%%%%%%%%%%%%

% Let's define the parameters of the $\ATTN$ layers for $h_f$ only attending to itself. 
% If we take as query and key matrix as

% \begin{align*}
%     & W_Q = W_K = \left(\begin{matrix}
%                     0      &\hdots &0      &0          \\
%                     \vdots &\ddots &\vdots &\vdots     \\
%                     0      &\hdots &0      &0          \\
%                     0      &\hdots &0      &\maxrep    \\
%                   \end{matrix}\right)
% \end{align*}

% \noindent then the keys and values of the vectors are 

% \begin{align*}
%     W_Q h_i = q_i = W_K h_i = k_i &= \begin{cases}
%                             (0, \dots, 0, \maxrep*1)     &\text{ if } i = f    \\
%                             (0, \dots, 0, \maxrep*-1)    &\text{ if } i \neq f \\
%                           \end{cases}
%                        &= \begin{cases}
%                             (0, \dots, 0, \maxrep)    &\text{ if } i = f    \\
%                             (0, \dots, 0, \minrep)    &\text{ if } i \neq f \\
%                           \end{cases}
% \end{align*}

% \noindent so the attention values for $h_f$ are $1$ for itself and 0 for the rest.

% \begin{align*}
%     s_f & = \softmax (&\langle q_f, k_1 \rangle,& &\ldots,& &\langle q_f, k_{f-1}\rangle,& &\langle q_f, k_f\rangle) &&\| (0,\ldots, 0) \\
%         & = \softmax (&\maxrep * \minrep,&        &\ldots,& &\maxrep * \minrep,          & &\maxrep*\maxrep)         &&\| (0,\ldots, 0) \\
%         & = \softmax (&          \minrep,&        &\ldots,& &          \minrep,          & &        \maxrep)         &&\| (0,\ldots, 0) \\
%         & =           &                0,&        &\ldots,& &                0,          & &               1,        && 0,\ldots, 0 \\
% \end{align*}

% \[
% (s_f)_i = \begin{cases}
%             1 &\text{ if } i = f    \\
%             0 &\text{ if } i \neq f \\
%           \end{cases}
% \]


%%%%%%%%%%%%%%%%%%%%%%%%%%%%%%%%%%%%%%%%%%%%%%%%%%
%%%%%%%%%%% END: ONLY ATTEND TO ITSELF %%%%%%%%%%%
%%%%%%%%%%%%%%%%%%%%%%%%%%%%%%%%%%%%%%%%%%%%%%%%%%

\medskip

%%%%%%%%%%%%%%%%%%%%%%%%%%%%%%%%%%%%%%%%%%%%
%%%%%%%%%%%% START: FIRST LAYER %%%%%%%%%%%%
%%%%%%%%%%%%%%%%%%%%%%%%%%%%%%%%%%%%%%%%%%%%

\santi{dividir claramente en primera y segunda layer. Podés usar itemize o paragrph.}
\rafa{estoy de acuerdo, pero no estoy seguro cómo hacerlo}


Recall from the definition of the $\ATTN$ mechanism (\ref{alg:def_attn}) that the output of the $\ATTN$ layer is $h'_i = W_O \sum_{j=1}^a (s_i)_j v_j$. 
By this point, $h'_f$ is looking like 

\[
h'_f = W_O \sum_{j=1}^a (s_f)_j v_j = W_O v_f = W_O (W_V \widetilde{h_f})
\]

\noindent Then, defining $W_O$ and $W_V$ as follows

\begin{align*}
    W_V &= id
    & 
    W_O &= \left(\begin{matrix}
                -id^{\Rdd}                 &0^{\Rdd}                   \\
                M                          &0^{\Rdd}                   
          \end{matrix}\right)
\end{align*}

\noindent makes us get $ h'_f = (-(h_f)_1, \dots, -(h_f)_d, (Mh_f)_1, \dots, (Mh_f)_d)$, which added to $\widetilde{h_f}$ gives us what we were looking for after the first layer.


%%%%%%%%%%%%%%%%%%%%%%%%%%%%%%%%%%%%%%%%%%
%%%%%%%%%%%% END: FIRST LAYER %%%%%%%%%%%%
%%%%%%%%%%%%%%%%%%%%%%%%%%%%%%%%%%%%%%%%%%

\medskip

%%%%%%%%%%%%%%%%%%%%%%%%%%%%%%%%%%%%%%%%%%%%%
%%%%%%%%%%%% START: SECOND LAYER %%%%%%%%%%%%
%%%%%%%%%%%%%%%%%%%%%%%%%%%%%%%%%%%%%%%%%%%%%

For the second $\ATTN$ layer, we use the same mechanism for making $\widetilde{h_f}$ only attend to itself.
And defining $W_V$ and $W_O$ as follows

\begin{align*}
    W_V &= id
    & 
    W_O &= \left(\begin{matrix}
                0^{\Rdd}                   &id^{\Rdd}                  \\
                0^{\Rdd}                   &-id^{\Rdd}                 \\
          \end{matrix}\right)
\end{align*}

\noindent we get $ h'_f = ((Mh_f)_1, \dots, -(Mh_f)_d, -(Mh_f)_1, \dots, -(Mh_f)_d)$, which added to $\widetilde{h_f}$ gives us what we were looking for after the second layer.


%%%%%%%%%%%%%%%%%%%%%%%%%%%%%%%%%%%%%%%%%%%
%%%%%%%%%%%% END: SECOND LAYER %%%%%%%%%%%%
%%%%%%%%%%%%%%%%%%%%%%%%%%%%%%%%%%%%%%%%%%%

\bigskip

Notice that after these two layers, the values of $h_i$ for all $i \neq f$ had been modified in ways not studied in this work.
\bigrafa{no sé cómo decir esto. El comentario anterior de que habían sido corrupted hacía referencia a esto.}

